\section{Discussion}\label{sec:discussion}

The reassessment changes the interpretation of the original three-stage result in two distinct ways. First, under a standard modern full-history SPNE definition, linear Hotelling price instability prevents a pure-strategy SPNE because at least one feasible price subgame has no pure Nash equilibrium. This observation is conceptually prior to any calculation at maximal locations, but it is not the paper's main novelty because the underlying price-instability problem is well known and was recognized in the literature surrounding the original model.

Second, and more importantly for the original backward-induction calculation, even a favorable on-path reading does not restore the reported quality-stage characterization. At the maximal locations, the regular pricing formula is valid only for quality gaps within its branch. A firm can make a finite quality deviation that changes the subsequent price equilibrium to exclusion. That global deviation overturns the symmetric candidate for $1/9<\lambda<4/27$, while the asymmetric exclusion profiles exist throughout the original interval $(1/9,2/9)$. Thus the central issue is not a small algebraic correction to a first-order condition but the use of a local continuation formula as if it described the whole strategy domain.

These results should not be read as a complete solution of every equilibrium-selection issue in the original game. We do not characterize mixed-strategy continuation equilibria at histories where pure price equilibria fail, and we do not claim a complete list of pure quality equilibria at maximal locations. The contribution is narrower: it identifies a certified global deviation and certified additional equilibria sufficient to invalidate the reported symmetric uniqueness characterization, while keeping the original noncooperative pricing structure unchanged.

The welfare results reinforce that distinction. Where the reported symmetric profile and the certified exclusion profiles coexist, the exclusion equilibria yield higher total welfare in the baseline model. This does not imply that exclusion is generally efficient or desirable. Conditional on maximal locations, the exclusion profile happens to implement the fixed-location planner optimum; the remaining loss relative to first best is spatial. A policy conclusion would require instruments and selection mechanisms absent from the model.
